\summarychapter{11}{Emmanuel!}{Christ in the Creeds}

\begin{inbrief}
After the Trinity came the question of Christ himself: how are the divine and the human related in the
one Savior? A century of debate between the schools of Alexandria and Antioch, sharpened by the rivalry
of the great bishops, produced three rejected answers and, at Chalcedon in 451, the classic definition
of one person in two natures.
\question{By what steps did the church arrive at its classic statement of the incarnation, and what does that statement mean?}
\end{inbrief}

Near Mount Hermon Jesus asked his disciples, ``Who do you say I am?'' and Peter answered, ``You are the
Messiah, the Son of the living God.'' People have given a thousand other answers---rabbi of love,
social revolutionary, misguided dreamer---but the church has confessed with Peter. Theologians call
the mystery the incarnation; hymns praise ``Emmanuel,'' God with us. Chalcedon declared Christ
``complete in Godhead and complete in manhood,'' in two natures forming one person. Most Roman
Catholic, Orthodox and Protestant Christians hold this as orthodox; some Christians in Egypt, Syria and
India do not.

\ssection{The Word Became Flesh}
Christology asks who Jesus Christ was. That the question exists at all is a monument to his
uniqueness: Islam has no ``Mohammedology.'' The imperial age did not invent the question; the God-man
was central to worship long before he was central to thought, and the heart of Christian experience,
as J.~S.~Whale said, lies in the church's praise rather than its creed. Still, God in the flesh is a
reality, not a sentiment. The debates did not aim to explain Christ, who fits no class, and the creeds deserve credit for leaving the mystery a mystery. The mind's revolt against them, said Charles
Williams, is immature; the romantic ``simple Jesus'' will not serve a faith that has always believed
earth and heaven meet in him.

Two schools emphasized different sides of ``the Word became flesh.'' Alexandria stressed that the
divine Word took on human flesh and risked letting the divinity overshadow the humanity. Its early
voice, Origen, coined the term ``God-man'' and built a Platonic mysticism of the Word; Gregory of Nyssa
later taught that the one divine person united two distinct natures whose attributes could be
exchanged. Antioch read Scripture historically, stressed the concrete man of the Gospels and the
saving value of his obedient will, and kept the two natures so distinct that they could seem two
beings, the Word dwelling in the man Jesus as in a temple; its risk was to let the humanity overshadow
the divinity.

Politics inflamed the debate. In the official church power gathered in a few hands: bishops of
provincial capitals became archbishops, and those of Rome, Constantinople, Alexandria and Antioch
were called patriarchs, each seeking to extend his see's prestige. Alexandria and Rome generally
supported each other out of resentment of the upstart ``new Rome''; Antioch, Alexandria's old rival,
sided with Constantinople.

\ssection{Major Heresies About Christ}
Between 350 and 450 each heresy forced a clearer answer. The first came from Apollinaris of Laodicea,
a younger friend of Athanasius. Reacting against Antioch, he held that at the incarnation the divine
Word took the place of the rational soul in a human body, giving ``one enfleshed nature of the divine
Word.'' Critics replied that the Gospels show a complete man and that a Christ without a human mind
could not redeem the whole person: ``What has not been assumed cannot be restored,'' said Gregory of
Nazianzus. The Council of Constantinople in 381 silenced the teaching.

The second is linked with Nestorius, a famous preacher of Antioch made bishop of Constantinople in
428. Following his teacher Theodore of Mopsuestia, he rejected the popular title of Mary as
``God-bearer, Mother of God,'' since God cannot be born or be three days old. He did not deny Christ's
deity, but he described the union of the natures as a moral joining of wills---``I hold the natures
apart, but unite the worship''---so that he seemed to teach two persons. He was condemned more for
politics than doctrine. Hated in the capital for his attacks on Jews and heretics, he made an enemy of
Cyril, patriarch of Alexandria (412--444), an able but ruthless theologian, who stirred up charges and
won over Pope Celestine. At the Council of Ephesus in 431, called by Theodosius II, Cyril had
Nestorius deposed before his Syrian supporters arrived; they, under John of Antioch, then condemned
Cyril; the Roman legates sided with Cyril. The historian Williston Walker thought few quarrels in the church's past so repellent. Nestorius was expelled and died in exile in Egypt about 450; how far he
was really heretical remains unclear. In his autobiography he said he only wished to stress that
Jesus was born fully human.

His followers fled to Persia and founded the Nestorian Church, with vigorous monasticism, theology
and missions reaching Malabar in India, Turkestan and, between 780 and 823, Tibet and central China,
until the persecution of Tamerlane (1380) destroyed the Asian mission. About eighty thousand members
remain in the Near East and India and twenty-five thousand in America. The church still uses
Nestorius's name for the failure to see the natures as truly united. Shelley notes that some liberal
views are called Nestorian, making Jesus differ from us only in degree, a model rather than a divine
savior, and that evangelicals fall into the same error when they treat Jesus as a self-help example.

\ssection{Establishing the Boundaries of Truth}
The third heresy followed quickly. Eutyches, head of a monastery near Constantinople, taught one
nature (monophysitism): Christ's humanity is absorbed in his deity like a drop of honey in the
sea---which endangered the whole doctrine of redemption. Flavian, patriarch of Constantinople, condemned
him, but Dioscorus of Alexandria, heir to Cyril's ideas, backed him and at a council in Ephesus in
449 had him restored. Pope Leo I called it the ``Robber Council'' and asked for a new one. The emperor
Marcian summoned it to Chalcedon, near Constantinople, in 451. Nearly four hundred bishops indicted
Dioscorus and, though reluctant to add to Nicaea, defined Christ as truly God and truly man, in two
natures ``without confusion, without change, without division, without separation,'' the properties
of each preserved and coming together in one person. Thus the church said against Arius that Jesus
is truly God, against Apollinaris that he is truly man, against Eutyches that the natures are not turned into anything else, and, in answer to Nestorius, that he is one person.

Many Christians in the Near East rejected Chalcedon and held to one nature. This, together with the
decline of Byzantine power in the outlying provinces, produced the Monophysite churches: the Coptic
Church of Egypt with its sister church in Ethiopia, and the Jacobite Church of Syria, whose members
are now mostly in South India. Chalcedon did not explain how deity and humanity unite; the event is
unique and resists explanation. Its merit, Shelley concludes, is the boundary it set: it built a
fence and said that the mystery of the God-man lies inside it. We may wish for clearer terms, but we
dare not say less.

\begin{keyfacts}
\keypeople{Origen; Gregory of Nyssa; Apollinaris; Gregory of Nazianzus; Theodore of Mopsuestia;
Nestorius; Cyril of Alexandria; Pope Celestine; Theodosius II; John of Antioch; Eutyches; Flavian;
Dioscorus; Pope Leo I; the emperor Marcian.}
\keydates{%
  \kd{381}{Council of Constantinople rejects Apollinaris's teaching}
  \kd{428}{Nestorius becomes bishop of Constantinople}
  \kd{431}{Council of Ephesus deposes Nestorius}
  \kd{449}{The ``Robber Council'' of Ephesus restores Eutyches}
  \kd{451}{Council of Chalcedon: one person in two natures}
  \kd{780--823}{Nestorian missionaries reach Tibet and central China}
  \kd{1380}{Tamerlane's persecution destroys the Nestorian mission in Asia}}
\keyterms{Incarnation; Christology; God-man; God-bearer (Mother of God); patriarch; see;
Apollinarianism; Nestorianism; monophysitism; Chalcedonian definition; Coptic and Jacobite churches.}
\end{keyfacts}
